\section{The Software}\label{sec:software}
\paragraph{Design goals} SurroSelect is a small Python package that works on plain arrays: a matrix \texttt{X} of simulation inputs (one row per run) and a matrix \texttt{Y} of output fields (one row per run, a common grid). It does not call the user's solver, so it can be used with any simulation code, commercial or not, as long as the runs can be exported to \texttt{.npz} files. It depends only on NumPy, SciPy and scikit-learn; a tuned neural network is an optional extra (PyTorch and Optuna). Fits are deterministic for a given seed, and a fitted selector can be saved and reloaded.
\paragraph{Structure} The package has five parts: the pool members (\texttt{models}: POD with a Gaussian process, POD with wide-range kernel ridge, POD with a ridge polynomial, full-field Gaussian processes, all with closed-form or marginal-likelihood hyper-parameter fitting); the selector (\texttt{select}: cross-validated selection with ties, the basis check, the learning curve and forecast, and the no-usable-model flag); the report writer (\texttt{report}); the cost-benefit arithmetic (\texttt{cost}); and the command-line interface (\texttt{cli}). The optional network (\texttt{nn}) is chosen by the same cross-validation as every other member.
\paragraph{Use} The command \texttt{surroselect fit data.npz \allowbreak--target 0.01 \allowbreak--solver-seconds 30 \allowbreak--queries 10000 \allowbreak--model model.pkl} prints a report with the cross-validation table, the tie set, the basis check, the learning curve with the forecast, a warning when no pool member is usable (best cross-validated error above 0.25), and a cost paragraph that states after how many evaluations the stand-in is cheaper than the solver; \texttt{surroselect predict} and \texttt{surroselect evaluate} apply and score the saved model on new files. The same functions are available from Python (\texttt{SurrogateSelector}, \texttt{make\_report}, \texttt{cost\_benefit}); \cref{lst:case} shows a report and Supplement Sec.~\ref{sec:S-tool} lists the interface.
\paragraph{Verification} \NToolTests{} unit tests cover the closed-form leave-one-out identity against brute force, exact recovery of a polynomial, selection and tie reporting, the basis-limit flag, inversion of the power law, the command-line round trip, numerical stability near interpolation, the alternative error metric, scalar outputs, the optional network, the warning and the cost arithmetic; the benchmark package has \NBenchTests{} further tests, including mesh-convergence and identity checks for every problem solved with \texttt{scikit-fem}. A continuous-integration workflow runs the tool tests on two Python versions. The changelog records that the selection logic is unchanged since the first version, so the evaluation of that version applies to the current one; what was added in later versions (the warning, the cost paragraph, command-line options, input validation, a note when the learning curve is not monotone, in version 0.3.0 the empirical interval for the forecast, in 0.4.1 the opt-in rational member and in 0.5.0 the opt-in pole-deflated model) was evaluated separately or changes no selection result (\cref{sec:confirm,sec:forecast}).
\begin{table}[t]\centering\caption{Software availability.}\label{tab:softavail}\small
\begin{tabularx}{\textwidth}{lX}\toprule
Name and version & SurroSelect 0.5.0\\
Developers and contact & P N Bhargav Teja (Bhargavteja.pn2023@vitstudent.ac.in)\\
First release & 2026\\
Language and requirements & Python $\ge$ 3.9; NumPy, SciPy, scikit-learn; optional PyTorch and Optuna; runs on a laptop (no GPU)\\
Availability & on request from the corresponding author (no public repository or archive DOI yet); the source is in the project directory \texttt{surroselect/}\\
Licence & MIT (copyright holder P N Bhargav Teja)\\
Cost & free\\
Documentation & README with an interface summary, changelog, example script (\texttt{examples/demo.py})\\
\bottomrule\end{tabularx}\end{table}
